\hypertarget{dsu-without-path-compression}{%
\subsection{DSU without Path
Compression}\label{dsu-without-path-compression}}

\begin{itemize}
\tightlist
\item
  When we are using DSU without path compression, in scenarios like DSU
  with rollbacks, we can need to keep the DSU time complexity to a
  minimum.
\item
  This is also acheived by a technique similar to small to large. Here,
  when merging two components, we set the parent of the smaller
  component to the larger component.
\item
  Each time we traverse an edge to a parent, the overall subtree size at
  that parent is doubled. Implies \texttt{O(logn)} for find operation.
\end{itemize}

\hypertarget{bit-tricks}{%
\subsection{BIT TRICKS}\label{bit-tricks}}

\begin{itemize}
\tightlist
\item
  To get the value of the least significant bit that is one
  \texttt{V\ =\ (\ x\ \ \&\ \ (-x)\ )}
\item
  For Subarray \texttt{\&} think about suffix subarray \texttt{\&} and
  traverse form right to left since number of distinct values for
  arr{[}L, r{]} is atmost \texttt{64}, number of set bits decreases
  monotonically for \texttt{r\ \textgreater{}=\ L} as \texttt{r}
  increases.
\end{itemize}

\hypertarget{math}{%
\subsection{MATH}\label{math}}

\begin{itemize}
\tightlist
\item
  \texttt{(x\ mod\ a)\ mod\ (ak)\ =\ (x\ mod\ (ak))\ mod\ (a)\ =\ x\ mod\ (a)}
\item
  If we want to determine the probability of moving from the state
  \texttt{X} to state \texttt{Y} over \texttt{K} steps - use Markov.
  \{matrix exponentiation\}
\item
  It is better to take \texttt{\%\ MOD} of the matrix before
  exponentiation or multiplication to reduce the number of modulo
  operations.
\end{itemize}

\hypertarget{graphs}{%
\subsection{GRAPHS}\label{graphs}}

\begin{itemize}
\item
  A DAG \texttt{G} has at least one vertex with in-degree \texttt{0} and
  one vertex with out-degree \texttt{0}
\item
  In the DFS tree of an ``undirected graph'' \textbf{edges can only go
  from a descendant to an ancestor} i.e, No cross edges. The back-edges
  of the graph all connect a vertex with its descendant in the spanning
  tree. This is why DFS tree is so useful.
\end{itemize}

\hypertarget{cmp}{%
\subsection{CMP}\label{cmp}}

\begin{itemize}
\tightlist
\item
  SET
\end{itemize}

\begin{verbatim}
 auto cmp = [&](ll a, ll b)
 {
    return a > b;
 }; 
 set<ll, decltype(cmp)> s(cmp); 
\end{verbatim}

\begin{itemize}
\tightlist
\item
  MULTISET
\end{itemize}

\begin{verbatim}
  auto cmp = [&](ll a, ll b)
  {
    return a > b;
  }; 
  multiset<ll, decltype(cmp)> s(cmp); 
\end{verbatim}

\begin{itemize}
\tightlist
\item
  PRIORITY QUEUE
\end{itemize}

\begin{verbatim}
  auto cmp = [&](ll a, ll b)
  {
    return a > b;
  }; 
  priority_queue<ll, vl, decltype(cmp)> pq(cmp); 
\end{verbatim}

\begin{verbatim}
  priority_queue<ll, vl, greater<ll>> pq; 
  // min heap, default max heap
\end{verbatim}

\begin{itemize}
\tightlist
\item
  ORDERED SET (NO fsanitize)
\end{itemize}

\begin{verbatim}
  auto cmp = [&](ll a, ll b)
  {
    return a > b;
  }; 
  ordered_set<ll, decltype(cmp)> s(cmp); 
\end{verbatim}

\begin{itemize}
\tightlist
\item
  USING STRUCT
\end{itemize}

\begin{verbatim}
struct student
{
  string s; 
  ll r;
  bool operator<(const student &other) const
  {
    return r < other.r;
  }
};
\end{verbatim}

\hypertarget{two-pointers}{%
\subsection{TWO POINTERS}\label{two-pointers}}

\begin{itemize}
\tightlist
\item
  Works if \texttt{l\textquotesingle{}\ \textgreater{}\ l} then
  \texttt{r\textquotesingle{}\ \textgreater{}=\ r}
\end{itemize}

\begin{verbatim}
// template for 2 pointers
ll head = -1, tail = 0;
while (tail < n)
{
  while ((head + 1 < n) && check())
  {
    head++;
    // upd
  }
  ans = max(ans, head - tail + 1);
  if (tail < head)
  {
    tail++;
    // upd
  }
  else
  {
    tail++;
    head = tail - 1;
    // reset
  }
}
\end{verbatim}
